\documentclass[12pt, %font size
               a4paper, %paper type
               twoside, % two sided printing
               openright, % start new chapter on right side only ( inserts blank pages )
               abstract=on, % Use an abstract
               DIV=11,      % This parameter organizes the borders (detailed explanation at http://texdoc.net/texmf-dist/doc/latex/koma-script/scrguide.pdf
               BCOR=8mm]{scrbook} % BCOR sets the space, used by the type of  book. (e.g. glued, hard cover..)
%scrreprt is used for larger texts with chapters (Master or Bachelor Thesis)
%scrartcl is used when there are no chapters ( for smaller paper)


\usepackage[latin1]{inputenc} % enables unicode support for windows
\usepackage[english]{babel} % sets up new english typography reform
\usepackage{csquotes} % babel demands it
\usepackage[unicode=true]{hyperref} % enables use of metadata for pdfs and hyperlinks within a document
\usepackage[natbib,maxnames=2,maxbibnames=100,style=authoryear-comp,uniquename=full,firstinits,doi=false,backend=biber,backref,hyperref]{biblatex} % advanced bibliography support
\usepackage[usenames,dvipsnames,hyperref]{xcolor} % enables more advanced color support for hyperref
\hypersetup{colorlinks=true, %flag for prints
                        hidelinks,  % this option would hide links for the print version of your thesis
			linkcolor=red!35!black,    %definition of the link color
			citecolor=green!35!black,  %definition of the cite color
			urlcolor=magenta!35!black, %definition of the url color
			pdfauthor=Tobias Hartmann, % Optional: Specify the author of the pdf
			pdftitle= Basic Latex template}      % Optional: Specify the title within the pdf
\usepackage{subfiles} %This package is used for subfiles
\usepackage{tabu}     % provides advanced tables
 \usepackage{array,multirow}
\bibliography{bibFiles/relatedWork.bib}% Include all the bibliography files
\usepackage{booktabs} % enables reference bookstyle tables
\usepackage[format=plain, labelfont=bf]{caption}
\usepackage[capitalize,noabbrev]{cleveref}
\usepackage{subcaption} % enables use multiple figures in a figure
\captionsetup{compatibility=false}
\usepackage{eurosym} %includes the euro symbol 
\usepackage{enumitem} % allows customization of enumeration and itemize environment
\usepackage{graphicx} % enables loading of graphics
\usepackage{tikz} % drawing vector graphics in latex
%\usepackage{parskip} %alternatively parskip replaces paragraph indentation by increased in-betweeen-paragraph linespacing 
\usepackage{setspace} % helps setup line spacing
%\onehalfspacing % increases linespacing to one and half
\usepackage{placeins} % provides FloatBarrier
%\usepackage[miktex]{gnuplottex} % gnuplot within latex. May be obsolete with pylab.
\usepackage[ruled,vlined]{algorithm2e} %algorithm package
\linespread{1.1} % Definition of the linespread
\usepackage[tbtags]{mathtools}
\DeclareMathOperator*{\argmax}{arg\,max} %argmax command

%tikz helps to draw nice pictures with a lot of effort for advanced users
\usetikzlibrary{positioning,shapes,shadows,arrows, backgrounds}
\usepackage{verbatim}
\usepackage{tikz-3dplot}


%some definitions for the cref package
\crefname{algocf}{Algorithm}{Algorithms}
\crefname{table}{Table}{Tables}
\crefname{chapter}{Chapter}{Chapters}
\crefname{equation}{Equation}{Equations}
\crefname{section}{Section}{Sections}


\tikzset{
    tri/.style={
        draw,
        shape border rotate=90,
        isosceles triangle,
        isosceles triangle apex angle=60,
        node distance=1cm,
        minimum height=4em
    }
}



\begin{document}
\frontmatter

\begin{titlepage}
\begin{center}

% Upper part of the page. The '~' is needed because \\
% only works if a paragraph has started.
\includegraphics{logo2}\\[1cm]

\textsc{\LARGE Rheinische\\[5mm] Friedrich-Wilhelms-Universität Bonn}\\[1.5cm]

\textsc{\Large Master thesis}\\[1.5cm]

% Title
{ \Large \bfseries Basic \LaTeX \, Template }\\[1.4cm]

% Author and supervisor
\begin{minipage}[t]{0.4\textwidth}
\begin{flushleft} \large
\emph{Author:}\\
Max \textsc{Mustermann}
\end{flushleft}
\end{minipage}
\begin{minipage}[t]{0.5\textwidth}
\begin{flushright} \large
\emph{First Examiner:} \\
Prof.~Dr.~given name \textsc{name} \\[0.5cm]
\emph{Second Examiner:} \\
Prof.~Dr.~given name.~\textsc{name} \\[0.5cm]

\emph{Advisor:} \\
given name \textsc{name} \\[0.5cm]
%\emph{Abteilung:} \\
%Autonome Intelligente Systeme
\end{flushright}
\end{minipage}

\vfill

% Bottom of the page
{\large Submitted:\hspace{1cm} \today}

\end{center}
\end{titlepage}

\pagestyle{headings}  % switches on printing of running heads
\title{Basic \LaTeX \, Template}
%\subtitle{Master thesis}
\author{Tobias Hartmann\\ \begin{minipage}{8cm}\centering \small Friedrich-Wilhelms-Universität Bonn\\ \small (group)\end{minipage}}

\vspace{4cm}

\cleardoublepage
\thispagestyle{empty}
{\noindent%
\huge{\textbf{\textsf{Declaration of Authorship}}}
}
\vspace{2cm}
\begin{flushleft}
\noindent%
I declare that the work presented here is original and the result of my own investigations.
Formulations and ideas taken from other sources are cited as such.

It has not been submitted, either in part or whole, for a degree at this or any other university.
\end{flushleft}

\vspace{8cm}
\noindent%
\rule[1em]{8em}{0.5pt}  \hfill \rule[1em]{8em}{0.5pt}\\ % This prints a line to write the date
Location, Date \hfill Signature\\


\cleardoublepage


\chapter*{Abstract}
\thispagestyle{empty}
Describe your approach and results within the abstract.


\newpage
\pagenumbering{roman}
\setcounter{page}{1}
\tableofcontents
\listoffigures
\listofalgorithms

\newpage
%\setcounter{page}{1}
%\pagenumbering{arabic}
\mainmatter
\subfile{subfiles/Introduction/Introduction}

\chapter{Methods you should probably use}
In general the DIV option should be used for specification of borders.
It allows a correct typography of the document.
The geometry package should be avoided whenever possible.

\section{Table of Contents}
Within the table of contents page numbering usually is in roman letters.
Afterwards the introduction should start with page number one in arabic letters.
The document class scrbook does this by default, when frontmatter, mainmatter and backmatter are defined.
Therefore, you should insert the \textbackslash frontmatter command before your table of contents, the \textbackslash mainmatter command before your introduction 
and finally the \textbackslash backmatter command before your appendix or bibliography.
Please have a look at the source code of this document, all packages are explained within the comments.


\section{Indentation}
Sometimes the indentation of latex at first feels odd.
Anyway if you have a look at some books, you will find that mostly a new paragraph is marked by indentation, instead of usage of an empty line.

Therefore, the indentation at the begin of this line is a desired effect.

\section{Referencing}
Use the cleverrev package to reference \cref{chap:Introduction}.
The package allows to specify the reference style globally.
Therefore, it helps to refer in a consistent way.
Otherwise you may quickly have both way of references e.g. Figure and figure, since it must be written manually.

\section{Citations}
There are two ways to cite:
\begin{itemize}
\item \citet{muster} showed that ...
\item This is currently a hot research topic \citep{muster}
\end{itemize}
In your bibliography, you should use strings, using the @string syntax.
Otherwise you might quickly end up with a lot of different names for the same conference, which causes misleading information.
(Please have a look at the bibliography file.)

\section{Enumerations}
For enumerations the package enumitem should be used since it allows for much more customization then the basic enumeration.
An example from the documentation is given below.


A ``normal'' enumerate list with changed counter.
\begin{enumerate}[label=\alph*)]
	\item Test 1
	\item Test 2
	\item Test 3
		\begin{enumerate}[label=\arabic*)]
			\item Test 3.1
			\item Test 3.2
			\item Test 3.3
				\begin{enumerate}[label=\Roman*)]
					\item Test 3.3.1
					\item Test 3.3.2
					\item Test 3.3.3
				\end{enumerate}
		\end{enumerate}
\end{enumerate}

Another list where the indentation was removed.
\begin{enumerate}[leftmargin=0em, label=\alph*)]
	\item Test 1
	\item Test 2
	\item Test 3
		\begin{enumerate}[label=\arabic*)]
			\item Test 3.1
			\item Test 3.2
			\item Test 3.3
				\begin{enumerate}[label=\Roman*)]
					\item Test 3.3.1
					\item Test 3.3.2
					\item Test 3.3.3
				\end{enumerate}
		\end{enumerate}
\end{enumerate}


\section{Figures}
An example figure of nyan cats is given in \cref{fig:cats}.
It is a basic example how to use use sub figures.
Please note that you should provide a short name for every figure in these braces: $[$ $]$.
The short name is displayed within the table of figures, where no detailed explanation is suitable.
If you do not define a short name, the whole text of the caption will be displayed.
The example algorithm below provides a bad example, where the short name should be specified but is not.
You can use the FloatBarrier command if you want to make sure that a figure is within a specific section.

In general there should be no forced line brake within a caption!

\begin{figure}
          \centering
        \begin{subfigure}[b]{0.3\textwidth}
                \includegraphics[width=\textwidth]{img/cat.jpeg}
                \caption{Nyan cat}
                \label{fig:cat1}
        \end{subfigure}%
        ~ %add desired spacing between images, e. g. ~, \quad, \qquad, \hfill etc.
          %(or a blank line to force the subfigure onto a new line)
        \begin{subfigure}[b]{0.3\textwidth}
                \includegraphics[width=\textwidth]{img/cat.jpeg}
                \caption{Nyan cat}
                \label{fig:cat2} 
        \end{subfigure}
        ~ %add desired spacing between images, e. g. ~, \quad, \qquad, \hfill etc.
          %(or a blank line to force the subfigure onto a new line)
        \begin{subfigure}[b]{0.3\textwidth}
                \includegraphics[width=\textwidth]{img/cat.jpeg}
                \caption{Nyan cat}
                \label{fig:cat3}  
        \end{subfigure}
	\caption[Nyan cats]{Pictures of Nyan cats}\label{fig:cats}
\end{figure}


\FloatBarrier

\section{Tables}
In general there should be no vertical lines within tables.
One package for tables is the booktabs package.
Many examples can be found in the web.
A basic example, taken from the booktabs documentation, is given in \cref{tab::ex}.

\begin{table}
\centering
	\begin{tabular}{lrr} 
	\toprule
	\multicolumn{2}{c}{Studium}\\ \cmidrule{1-2}
	Fach & Dauer & Einkommen (\euro)\\ 
	\midrule 
	Info & 2 & 12,75 \\ \addlinespace
	MST & 6 & 8,20 \\ \addlinespace
	VWL & 14 & 10,00\\ 
	\bottomrule
\end{tabular}
\caption[Table Example]{A basic example from the booktabs package.}
\label{tab::ex}
\end{table}
\FloatBarrier

\subsection{The tabu package}
The tabu package provides some useful tools for advanced tables.



\section{Algorithms}
There are several packages for algorithms.
A nice one is the algorithm 2e package.
In \cref{alg:exp} a simple example from the algorithm2e documentation is given.


\begin{algorithm}[H]
	\SetAlgoLined
	\KwData{this text}
	\KwResult{how to write algorithm with \LaTeX2e }
	initialization\;
	\While{not at end of this document}{
		read current\;
		\eIf{understand}{
			go to next section\;
			current section becomes this one\;
		}{
			go back to the beginning of current section\;
		}
	}
	\caption{How to write algorithms (Small example from the algorithm2e documentation)}
	\label{alg:exp}
\end{algorithm}
\FloatBarrier

\section{Compiling The Document}
You have to compile the document in three steps:
\begin{itemize}
	\item Compile the document using xelatex: xelatex Template.tex
 \item Run biber to set up bibliography: biber Template
 \item Run xelatex again and get the final document: Xelatex Template.tex
\end{itemize}

\backmatter

\begin{singlespacing}
\printbibliography
\end{singlespacing}
\FloatBarrier

\end{document}

